\chapter{Разработка и реализация программного комплекса моделирования HeatSim3D}

\section{Архитектура и структура программного комплекса HeatSim3D}

Программный комплекс HeatSim3D предназначен для численного моделирования 
нестационарной теплопроводности в трёхмерных сплошных средах. Комплекс 
реализован в виде двух взаимодействующих компонентов: расчётного модуля, 
написанного на языке C++~\cite{stroustrup1986overview}, и графического интерфейса пользователя, 
реализованного на языке Python~\cite{python2021python}.

Такая архитектура обусловлена различными требованиями к компонентам 
системы. Расчётный модуль выполняет вычислительно интенсивные операции 
и требует высокой производительности, которую обеспечивает компилируемый 
язык C++. Графический интерфейс, напротив, не предъявляет высоких 
требований к производительности, однако требует гибкости при разработке 
средств визуализации, для чего язык Python предоставляет широкие 
возможности в виде библиотек matplotlib~\cite{bisong2019matplotlib} и pandas~\cite{robitaille2013astropy}.

\subsection{Общая структура и взаимодействие компонентов}

Общая структура программного комплекса включает два основных компонента. 
Расчётный модуль реализован в виде исполняемого файла heat3d.exe и 
выполняет непосредственное численное моделирование на основе явной 
конечно-разностной схемы FTCS. Графический интерфейс реализован в виде 
скрипта gui.py и обеспечивает взаимодействие пользователя с программой, 
а также визуализацию результатов расчёта.

Взаимодействие компонентов организовано следующим образом. Пользователь 
задаёт параметры расчёта через графический интерфейс, после чего 
интерфейс формирует команду запуска расчётного модуля с передачей 
параметров через аргументы командной строки. По завершении расчёта 
модуль сохраняет результаты в файлы формата CSV, которые графический 
интерфейс считывает и отображает в виде графиков и тепловых карт.

\subsection{Передача параметров через аргументы командной строки}

Все параметры расчёта передаются расчётному модулю через аргументы 
командной строки в формате ключ-значение. Такой подход обеспечивает 
полную независимость компонентов и позволяет запускать расчётный модуль 
как из графического интерфейса, так и непосредственно из командной строки.

Расчётный модуль поддерживает следующие параметры: количество узлов 
сетки по каждой из осей ($N_x$, $N_y$, $N_z$), коэффициент 
температуропроводности материала $a$, время моделирования $t_{end}$, 
начальная температура $T_0$, температуры на каждой из шести граней 
расчётной области ($T_{x^-}$, $T_{x^+}$, $T_{y^-}$, $T_{y^+}$, 
$T_{z^-}$, $T_{z^+}$), а также частота сохранения результатов. 
При отсутствии какого-либо параметра используется значение по умолчанию, 
что обеспечивает работоспособность модуля без обязательного указания 
всех параметров.

\section{Расчётный модуль решения уравнения теплопроводности}

Расчётный модуль реализован на языке C++ и представляет собой 
консольное приложение, выполняющее численное моделирование 
нестационарной теплопроводности на основе явной конечно-разностной 
схемы FTCS. Модуль компилируется в исполняемый файл heat3d.exe с 
использованием стандарта C++17.

\subsection{Структуры данных и организация трёхмерного массива}

Температурное поле представлено в виде одномерного массива типа 
\texttt{std::vector<double>}, хранящего значения температуры во всех 
узлах трёхмерной сетки. Преобразование трёхмерных индексов $(i, j, k)$ 
в линейный индекс одномерного массива выполняется по формуле:
\begin{equation}
    \text{idx}(i, j, k) = i \cdot N_y \cdot N_z + j \cdot N_z + k.
\end{equation}

Использование одномерного массива вместо трёхмерного обусловлено 
требованиями к производительности: последовательное расположение 
элементов в памяти обеспечивает эффективное кэширование данных 
процессором~\cite{кудухов2014особенности} и снижает накладные расходы при обращении к элементам 
массива.

Для хранения расчётных параметров используется структура \texttt{Params}, 
содержащая все входные параметры модели: размеры сетки, физические 
свойства материала, начальные и граничные условия, а также параметры 
вывода результатов. Значения по умолчанию задаются непосредственно 
в определении структуры, что обеспечивает корректную работу модуля 
при неполном наборе входных параметров.

\subsection{Реализация алгоритма FTCS и граничных условий}

Основной вычислительный цикл реализует последовательное применение 
явной разностной схемы FTCS. На каждом временном шаге для всех 
внутренних узлов сетки вычисляется новое значение температуры по 
формуле:
\begin{equation}
    T_{i,j,k}^{n+1} = T_{i,j,k}^{n} + 
    r_x \delta_x^2 T + 
    r_y \delta_y^2 T + 
    r_z \delta_z^2 T.
\end{equation}


Граничные условия Дирихле~\cite{дирихлерешении} применяются после каждого шага по времени 
путём присвоения заданных температур узлам, расположенным на гранях 
расчётной области. Для каждой из шести граней температура задаётся 
независимо, что позволяет моделировать произвольные условия нагрева.

\subsection{Автоматический выбор шага по времени}

Временной шаг $\Delta t$ вычисляется автоматически перед началом 
расчёта на основе условия устойчивости схемы FTCS с коэффициентом 
запаса $0{,}4$:
\begin{equation}
    \Delta t = \frac{0{,}4}{2a \left( 
    \dfrac{1}{\Delta x^2} + 
    \dfrac{1}{\Delta y^2} + 
    \dfrac{1}{\Delta z^2} \right)}.
\end{equation}

Автоматический выбор шага исключает возможность задания пользователем 
неустойчивого режима расчёта и обеспечивает корректное численное 
решение при любых допустимых параметрах сетки и материала.

\subsection{Параллелизация вычислений с использованием OpenMP}

Для повышения производительности расчётного модуля реализована 
параллелизация основного вычислительного цикла с использованием 
технологии OpenMP~\cite{chandra2001parallel}. Директива параллельного выполнения применяется 
к внешним циклам по пространственным координатам, что позволяет 
распределить вычисления по доступным ядрам процессора.

Применение директивы collapse(2) объединяет два внешних 
цикла в единое параллельное пространство итераций, что обеспечивает 
более равномерное распределение нагрузки между потоками при различных 
соотношениях размеров сетки. В результате скорость расчёта возрастает 
пропорционально числу доступных вычислительных ядер.

\section{Графический интерфейс пользователя программного комплекса HeatSim3D}

Графический интерфейс пользователя реализован на языке Python с 
использованием библиотеки tkinter и обеспечивает взаимодействие 
пользователя с программным комплексом HeatSim3D. Интерфейс позволяет 
задавать параметры расчёта, запускать и останавливать симуляцию, 
отслеживать ход выполнения расчёта, а также просматривать результаты 
в виде графиков и тепловых карт непосредственно в окне программы.

\subsection{Панель управления параметрами и настройками численной симуляции}

Панель управления расположена в левой части главного окна программы 
и содержит все необходимые элементы для настройки расчёта. 
Пользователю доступны следующие параметры: выбор материала из 
встроенной библиотеки, включающей девять твёрдых материалов с 
заданными теплофизическими свойствами; параметры расчётной сетки 
$N_x$, $N_y$, $N_z$; коэффициент температуропроводности $a$; 
время моделирования $t_{end}$; начальная температура $T_0$; 
температуры на каждой из шести граней расчётной области; 
частота сохранения результатов; выбор цветовой палитры для 
отображения тепловых карт.

При выборе материала из списка значение коэффициента 
температуропроводности подставляется автоматически, однако 
пользователь может изменить его вручную для моделирования 
материалов с произвольными свойствами.

В нижней части панели управления расположены индикатор хода 
выполнения расчёта, отображающий текущий шаг и температуру 
в центре области, кнопка запуска симуляции и кнопка остановки. 
Ниже расположена область журнала, в которой отображается вывод 
расчётного модуля в режиме реального времени.

\subsection{Визуализация результатов --- тепловые карты и срезы}

Результаты расчёта отображаются на пяти вкладках главного окна 
программы. Вкладка динамики температуры содержит графики изменения 
температуры в центре расчётной области и средней температуры по 
объёму в зависимости от времени. На каждом графике автоматически 
отмечаются минимальное и максимальное значения температуры с 
указанием их численных значений.

Вкладка срезов и теплового потока содержит четыре панели: тепловую 
карту среза в плоскости $z = L_z/2$, тепловую карту среза в 
плоскости $y = L_y/2$, визуализацию вектора теплового потока и 
профиль температуры вдоль центральной оси. Тепловые карты 
отображают распределение температуры в цветовой шкале, 
выбранной пользователем.

Вкладка трёхмерной визуализации содержит изолинии температурного 
поля в плоскости $z = L_z/2$ и трёхмерную поверхность 
$T(x, y)$ для финального шага расчёта. Вкладка профилей 
содержит графики распределения температуры вдоль осей $x$ и $y$ 
для начального, промежуточного и конечного шагов расчёта.

\subsection{Визуализация вектора теплового потока в расчётной области
}

Вектор теплового потока вычисляется на основе закона Фурье 
непосредственно из данных температурного поля, сохранённых 
расчётным модулем. Компоненты вектора теплового потока 
определяются как:
\begin{equation}
    \mathbf{q} = -\nabla T = 
    \left( -\frac{\partial T}{\partial x}, 
           -\frac{\partial T}{\partial y} \right).
\end{equation}

Градиент температуры~\cite{махоньков2013влияние} вычисляется численно с использованием 
центральных разностей средствами библиотеки numpy. Направление 
стрелок соответствует направлению теплового потока --- от 
области с более высокой температурой к области с более 
низкой температурой. Для наглядности отображения длина 
всех стрелок нормируется к единице.

\subsection{Сравнение результатов двух независимых расчётов между собой}

Программный комплекс предоставляет возможность сравнения 
результатов двух независимых расчётов. После выполнения 
первого расчёта пользователь сохраняет его результаты 
под меткой A, затем изменяет параметры и выполняет второй 
расчёт, результаты которого сохраняются под меткой B. 
На вкладке сравнения отображаются совмещённые графики 
динамики температуры для обоих расчётов, а также тепловые 
карты финального поля температуры, что позволяет наглядно 
оценить влияние изменённых параметров на результаты 
моделирования.

\subsection{Сборка в исполняемый файл и инсталлятор}


Для обеспечения возможности использования программного 
комплекса без установки среды Python графический интерфейс 
упаковывается в автономный исполняемый файл с помощью 
инструмента <<PyInstaller>>. Итоговый дистрибутив включает 
исполняемый файл графического интерфейса <<HeatSim3D.exe>> 
и расчётный модуль <<heat3d.exe>>.

Для удобства установки программного комплекса на целевом 
компьютере разработан инсталлятор <<HeatSim3D\_Setup.exe>>, 
созданный с использованием инструмента <<Inno Setup>>. 
Инсталлятор выполняет установку программы в директорию 
Program Files, создание ярлыка на рабочем столе и 
регистрацию программы в системном списке установленных 
приложений, что обеспечивает возможность корректного 
удаления программы стандартными средствами операционной 
системы.
